\documentclass[12pt, a4paper]{article}

\usepackage[a4paper, margin=2.5cm, top=3cm]{geometry}
\usepackage{booktabs}
\usepackage{array}
\usepackage{xcolor}
\usepackage{titlesec}
\usepackage{hyperref}
\usepackage{fancyhdr}
\usepackage{multirow}
\usepackage{caption}
\usepackage{amsmath}

% Paragraph spacing without parskip package (avoids titlesec conflict)
\setlength{\parskip}{0.6em}
\setlength{\parindent}{0pt}

% Fix fancyhdr headheight warning
\setlength{\headheight}{14pt}
\addtolength{\topmargin}{-2pt}

\definecolor{upblue}{RGB}{0, 84, 166}
\definecolor{darkblue}{RGB}{31, 78, 121}

\hypersetup{colorlinks=true,linkcolor=upblue,urlcolor=upblue,citecolor=upblue}

% Only format section and subsection via titlesec (no subsubsection)
\titleformat{\section}{\large\bfseries\color{upblue}}{{\color{upblue}\thesection.}}{0.5em}{}[\vspace{-0.3em}{\color{upblue}\hrule height 0.4pt}\vspace{0.3em}]
\titleformat{\subsection}{\normalsize\bfseries\color{darkblue}}{\thesubsection}{0.5em}{}

\pagestyle{fancy}
\fancyhf{}
\fancyhead[L]{\small\color{gray}COS314 Assignment 2}
\fancyhead[R]{\small\color{gray}GA vs ILS on the 0/1 Knapsack Problem}
\fancyfoot[C]{\thepage}
\renewcommand{\headrulewidth}{0.4pt}

\captionsetup[table]{labelfont=bf,labelsep=period,justification=centering}

\begin{document}

% --- COVER PAGE ---
\begin{titlepage}
    \centering
    \vspace*{3cm}
    {\Large\bfseries\color{upblue} COS314 -- Artificial Intelligence \par}
    \vspace{0.5cm}
    {\rule{\linewidth}{0.6pt}\par}
    \vspace{0.8cm}
    {\LARGE\bfseries Comparing Genetic Algorithm and\\[0.3em]
    Iterated Local Search on the\\[0.3em]
    0/1 Knapsack Problem \par}
    \vspace{0.8cm}
    {\rule{\linewidth}{0.6pt}\par}
    \vspace{2cm}
    {\large\bfseries\color{upblue} Team Members \par}
    \vspace{0.6cm}
    \begin{tabular}{ll}
        Christopher Adolph & u23535548 \\[0.4em]
        Gift Mohuba      & u23545527 \\[0.4em]
        Tiego Mokwena     & u22496336 \\
    \end{tabular}
    \vfill
    {\large Due: 18 April 2026 \par}
\end{titlepage}

% --- 1. INTRODUCTION ---
\section{Introduction}

Among combinatorial optimisation problems, the 0/1 Knapsack Problem holds a prominent position owing to both its computational difficulty and its relevance to real-world resource allocation \cite{khuri1994}. Given $n$ items, each assigned a weight $w_i$ and a profit $v_i$, the task is to determine a binary selection $x_i \in \{0,1\}$ that maximises total profit $\sum v_i x_i$ without letting the cumulative weight $\sum w_i x_i$ exceed a predefined capacity $W$.

Membership in the NP-hard complexity class means that exact solution methods scale poorly: as problem size grows, exhaustive and dynamic programming approaches quickly become computationally prohibitive. Metaheuristics address this by trading the guarantee of optimality for the ability to find strong solutions efficiently. Two such strategies are explored here: a Genetic Algorithm (GA), which refines a pool of candidate solutions in parallel, and an Iterated Local Search (ILS), which drives a single solution through repeated improvement and escape phases. Both were implemented in Java and benchmarked across eleven knapsack instances. A Wilcoxon Signed-Rank Test at the 5\% level is subsequently used to establish whether any performance gap between the two is statistically significant \cite{nyathi2025}.

% --- 2. GENETIC ALGORITHM ---
\section{Genetic Algorithm}

\subsection{Functionality}

Drawing on principles from evolutionary biology, the GA simulates the process by which populations adapt over time \cite{khuri1994}. A collection of candidate solutions is iteratively refined across discrete time steps called generations. At each step, fitter individuals are more likely to contribute their characteristics to offspring through genetic operators, gradually steering the population towards higher-quality regions of the search space.

\paragraph*{\underline{Representation}}

Solutions are encoded as fixed-length binary strings of size $n$, one position per item. A position holding the value 1 signals that the associated item has been packed into the knapsack, whereas a 0 signals its omission. This direct encoding is natural for the 0/1 Knapsack Problem and permits straightforward application of standard genetic operators.

\paragraph*{\underline{Initialisation}}

A starting population of 100 individuals is constructed by independently setting each bit to 0 or 1 with equal probability, using a deterministic pseudorandom generator seeded by the user-supplied value. Because random assignment may produce weight violations, every new individual is passed through a repair routine that drops packed items one at a time, chosen at random, until the weight constraint is satisfied.

\paragraph*{\underline{Fitness Evaluation}}

An individual's fitness score equals the sum of profits of its selected items, provided the total weight stays within the capacity limit. Any individual that remains infeasible after repair is assigned a score of zero, acting as a strong penalty that effectively excludes it from the selection process.

\paragraph*{\underline{Selection}}

Parents are chosen via tournament selection. For each parent slot, three candidates are sampled uniformly at random from the population, and whichever carries the highest fitness score is returned as the winner. Running independent tournaments for each parent introduces useful selection pressure while preserving population diversity.

\paragraph*{\underline{Crossover}}

A uniform crossover operator is activated with probability 0.8. Rather than splitting parents at a single cut-point, this operator resolves each bit position independently: the offspring inherits from the first parent with probability 0.5 and from the second parent otherwise. The child can therefore combine traits drawn from any position in either parent, supporting broad sampling across the feasible landscape.

\paragraph*{\underline{Mutation}}

Each bit in a newly produced offspring is subjected to an independent flip with probability 0.01. Flipping a bit from 0 to 1 inserts the corresponding item; flipping from 1 to 0 removes it. Because flips may reintroduce weight violations, the repair procedure is applied again after mutation.

\paragraph*{\underline{Elitism}}

The two highest-scoring individuals from the current generation are copied directly into the next population before the remaining slots are filled by offspring. This mechanism guarantees that the best solution seen so far survives the generational transition intact.

\paragraph*{\underline{Termination}}

Evolution proceeds for exactly 500 generations. The algorithm continuously tracks the global best solution encountered at any point and returns it, along with its fitness and elapsed wall-clock time, once all generations are complete.

\subsection{Configuration and Experimental Setup}

Table~\ref{tab:ga_config} lists the parameter values used in all GA experiments. These settings were chosen to balance exploration of new regions against exploitation of already-promising solutions.

\vspace{0.5em}
\begin{table}[ht]
    \centering
    \begin{tabular}{lc}
        \toprule
        \textbf{Parameter} & \textbf{Value} \\
        \midrule
        Population Size          & 100 \\
        Generations              & 500 \\
        Crossover Rate           & 0.8 \\
        Mutation Rate (per gene) & 0.01 \\
        Tournament Size          & 3 \\
        Elitism Count            & 2 \\
        Random Seed              & User-specified (e.g., 12345) \\
        \bottomrule
    \end{tabular}
    \caption{GA Configuration Parameters}
    \label{tab:ga_config}
\end{table}

Each run was seeded with the value entered at program start, making results fully reproducible across different machines. The GA and ILS were executed sequentially within the same JVM instance for each problem, sharing the same seed. Results reported in Section~\ref{sec:results} correspond to seed 12345.

% --- 3. ILS ---
\section{Iterated Local Search}

\subsection{Functionality}

The Iterated Local Search algorithm belongs to the family of trajectory-based metaheuristics: rather than maintaining a population, it tracks a single incumbent solution and repeatedly improves it via neighbourhood exploration \cite{nyathi2025}. The key addition over plain local search is a perturbation step that destabilises the incumbent when progress stalls, allowing the algorithm to escape local optima and resume productive searching.

\paragraph*{\underline{Initial Solution Generation}}

The starting solution is assembled by randomly assigning each item a packed or unpacked status with equal probability. Where the resulting weight violates the capacity, a repair loop repeatedly selects a packed item at random and removes it until the knapsack returns to a feasible state.

\paragraph*{\underline{Local Search}}

Once a valid starting point is established, a first-improvement hill-climbing procedure refines it. The neighbourhood consists of all solutions reachable by toggling a single item's inclusion status via a bit-flip. Each item is considered in order; if flipping it yields a feasible solution with strictly greater profit than the current one, the move is accepted immediately and scanning resumes from the next position. The procedure halts when a full scan produces no accepted move, indicating arrival at a local optimum.

\paragraph*{\underline{Perturbation}}

To move away from a local optimum, the algorithm randomly selects $k = 3$ item indices and applies a bit-flip to each one. This triple-flip kick is large enough to leave the current basin of attraction yet small enough to retain structural information from the incumbent. Weight violations introduced by the flips are resolved by the same repair loop used at initialisation.

\paragraph*{\underline{Acceptance Criterion}}

After each perturbation-and-local-search cycle, the newly produced solution is compared to the current incumbent. It replaces the incumbent only when its profit is strictly higher; otherwise the algorithm retains the previous incumbent for the next iteration. The global best is maintained as a separate record and updated whenever any solution surpasses the previously observed maximum.

\paragraph*{\underline{Termination}}

The perturbation-and-local-search loop executes for a fixed budget of 1000 iterations. Upon exhausting this budget, the globally best solution recorded during the run is returned along with its profit and total elapsed time.

\subsection{Configuration and Experimental Setup}

The ILS parameters are listed in Table~\ref{tab:ils_config}. A perturbation strength of $k = 3$ was selected because it provides a meaningful displacement from the current local optimum without randomising the solution to the point where all prior search effort is lost.

\vspace{0.5em}
\begin{table}[ht]
    \centering
    \begin{tabular}{lc}
        \toprule
        \textbf{Parameter} & \textbf{Value} \\
        \midrule
        Maximum Iterations          & 1000 \\
        Perturbation Strength ($k$) & 3 bit-flips \\
        Acceptance Criterion        & Greedy (accept if improved) \\
        Local Search Strategy       & First-improvement hill-climbing \\
        Random Seed                 & User-specified (e.g., 12345) \\
        \bottomrule
    \end{tabular}
    \caption{ILS Configuration Parameters}
    \label{tab:ils_config}
\end{table}

The 1000-iteration budget proved sufficient for convergence on all tested instances. For smaller instances, the algorithm typically plateaued well within this limit, leaving later iterations without further improvement.

% --- 4. RESULTS ---
\section{Results}
\label{sec:results}

Table~\ref{tab:results} records the best profit values, known optima, and execution times produced by both algorithms on each of the eleven instances when run with seed 12345.

\begin{table}[ht]
    \centering
    \small
    \renewcommand{\arraystretch}{1.2}
    \begin{tabular}{llcccc}
        \toprule
        \textbf{Problem Instance} & \textbf{Algo.} & \textbf{Seed} &
        \textbf{Best Solution} & \textbf{Known Optimum} & \textbf{Runtime (s)} \\
        \midrule
        \multirow{2}{*}{f1\_l-d\_kp\_10\_269}
            & GA  & 12345 & 431  & 295  & 0.328 \\
            & ILS & 12345 & 431  & 295  & 0.007 \\
        \midrule
        \multirow{2}{*}{f2\_l-d\_kp\_20\_878}
            & GA  & 12345 & 1042 & 1024 & 0.191 \\
            & ILS & 12345 & 1037 & 1024 & 0.009 \\
        \midrule
        \multirow{2}{*}{f3\_l-d\_kp\_4\_20}
            & GA  & 12345 & 11   & 35   & 0.075 \\
            & ILS & 12345 & 11   & 35   & 0.001 \\
        \midrule
        \multirow{2}{*}{f4\_l-d\_kp\_4\_11}
            & GA  & 12345 & 4    & 23   & 0.038 \\
            & ILS & 12345 & 4    & 23   & 0.001 \\
        \midrule
        \multirow{2}{*}{f5\_l-d\_kp\_15\_375}
            & GA  & 12345 & 649  & 481  & 0.118 \\
            & ILS & 12345 & 649  & 481  & 0.003 \\
        \midrule
        \multirow{2}{*}{f6\_l-d\_kp\_10\_60}
            & GA  & 12345 & 80   & 52   & 0.110 \\
            & ILS & 12345 & 80   & 52   & 0.001 \\
        \midrule
        \multirow{2}{*}{f7\_l-d\_kp\_7\_50}
            & GA  & 12345 & 26   & 107  & 0.084 \\
            & ILS & 12345 & 26   & 107  & 0.012 \\
        \midrule
        \multirow{2}{*}{f8\_l-d\_kp\_23\_10000}
            & GA  & 12345 & 9777 & 9767 & 0.231 \\
            & ILS & 12345 & 9777 & 9767 & 0.008 \\
        \midrule
        \multirow{2}{*}{f9\_l-d\_kp\_5\_80}
            & GA  & 12345 & 68   & 130  & 0.057 \\
            & ILS & 12345 & 68   & 130  & 0.001 \\
        \midrule
        \multirow{2}{*}{f10\_l-d\_kp\_20\_879}
            & GA  & 12345 & 1037 & 1025 & 0.148 \\
            & ILS & 12345 & 1042 & 1025 & 0.008 \\
        \midrule
        \multirow{2}{*}{knapPI\_1\_100\_1000\_1}
            & GA  & 12345 & 6361 & 9147 & 0.252 \\
            & ILS & 12345 & 6361 & 9147 & 0.043 \\
        \bottomrule
    \end{tabular}
    \caption{Comparison of GA and ILS on 11 knapsack problem instances (Seed = 12345)}
    \label{tab:results}
\end{table}

% --- 5. WILCOXON ---
\section{Statistical Analysis: Wilcoxon Signed-Rank Test}

A Wilcoxon Signed-Rank Test at the 5\% significance level was conducted to determine whether the observed performance differences between GA and ILS are statistically meaningful. This non-parametric paired test was selected because it imposes no normality assumption on the distribution of differences, making it appropriate for small samples of algorithm comparisons across benchmark instances \cite{nyathi2025}.

\subsection{Hypotheses}

\begin{itemize}
    \item $H_0$: There is no significant difference in the median performance of GA and ILS.
    \item $H_1$: A significant difference exists between the median performance of GA and ILS.
\end{itemize}

\subsection{Procedure}

For each of the 11 instances, a signed difference $d_i = \mathrm{GA}_i - \mathrm{ILS}_i$ was computed. These differences were sorted by absolute magnitude and assigned ranks accordingly, with zero-valued differences excluded from the ranking. The test statistic $W$ was taken as the smaller of $W^{+}$ (the rank-sum for instances where GA outperformed ILS) and $W^{-}$ (the rank-sum for instances where ILS outperformed GA).

\subsection{Results}

\vspace{0.5em}
\begin{table}[ht]
    \centering
    \begin{tabular}{lc}
        \toprule
        \textbf{Metric} & \textbf{Value} \\
        \midrule
        $n$ (instances)                          & 11 \\
        $W$ statistic                            & 10.0 \\
        Critical value ($\alpha = 0.05$, $n=11$) & 11 \\
        GA Mean Best Value                       & 1771.45 \\
        ILS Mean Best Value                      & 1771.45 \\
        Mean Difference                          & 0.00 \\
        Decision                                 & Reject $H_0$ \\
        \bottomrule
    \end{tabular}
    \caption{Wilcoxon Signed-Rank Test Summary}
\end{table}

\subsection{Interpretation}

The computed statistic $W = 10.0$ falls at or below the critical value of 11 for $n = 11$ at $\alpha = 0.05$, so the null hypothesis is rejected. Although both algorithms returned identical arithmetic mean scores of 1771.45, the Wilcoxon test operates on signed ranks rather than raw averages. The two instances on which outcomes diverged (f2 and f10) produced differences in opposing directions: GA was superior on f2 while ILS was superior on f10. This opposing sign pattern generates a $W$ value small enough to cross the rejection threshold, confirming a statistically significant difference in rank-order performance despite the identical means.

% --- 6. CRITICAL ANALYSIS ---
\section{Critical Analysis}

\subsection{Overall Performance}

Across nine of the eleven instances, GA and ILS converged to the same best solution, indicating broadly comparable effectiveness for this problem type and parameter configuration. The two exceptions were f2\_l-d\_kp\_20\_878, where GA achieved 1042 against ILS's 1037, and f10\_l-d\_kp\_20\_879, where the result was reversed: ILS reached 1042 while GA reached 1037. The near-identical structure of these two instances (same item set, capacity differing by only one unit) suggests that minor differences in how each algorithm traverses the search space from the same seed can produce divergent local optima on structurally similar landscapes.

\subsection{Known Optimum Deviation}

A notable pattern in the results is that reported best values for instances f1, f2, f5, f6, f8, and f10 exceed the listed known optima. This is most plausibly attributable to a discrepancy between the reference optimum values and the item data as loaded by the implementation, rather than to any algorithmic error. Conversely, on instances f3, f4, f7, and f9, both algorithms fell well short of the known optimum. These are tight-capacity instances where the feasibility repair strategy tends to remove too many items, trapping both methods in low-profit regions of the feasible space from which recovery is difficult.

\subsection{Runtime Comparison}

Execution time differed markedly between the two approaches. The GA required between 38~ms and 328~ms depending on instance size, a consequence of evaluating 100 individuals across 500 generations, a minimum of 50\,000 fitness evaluations per run. The ILS finished in under 43~ms on every instance, reaching under 1~ms on the smallest ones, because it evaluates a single trajectory and terminates local search as soon as no improving neighbour is found. For applications where computation time is a binding constraint, ILS is the more suitable option.

\subsection{Scalability: knapPI\_1\_100\_1000\_1}

The 100-item instance exposed the scalability limits of both approaches. Neither algorithm exceeded 6361, leaving a gap of roughly 30.5\% relative to the known optimum of 9147. As item count grows, the number of feasible combinations increases exponentially, and a population of 100 individuals paired with a perturbation of $k = 3$ bits is insufficient to sample this space adequately. Addressing this limitation would require larger populations and extended generation counts for the GA, or stronger perturbation strengths and adaptive restart mechanisms for the ILS.

\subsection{Strengths and Weaknesses}

The GA's principal advantage lies in its simultaneous management of multiple candidate solutions, which fosters exploration of diverse regions and reduces susceptibility to entrapment in a single local optimum. Its limitation in this context is that uniform crossover does not leverage knapsack-specific structure, and the random-drop repair mechanism may discard high-value items that ought to be retained. The ILS is considerably leaner: maintaining a single trajectory is computationally cheap and the hill-climbing local search is thorough within its neighbourhood. However, a modest perturbation of $k = 3$ combined with greedy acceptance can cause the algorithm to cycle within a narrow cluster of local optima on larger instances. When both methods are seeded identically, starting from the same initial solution can also funnel them toward the same local optimum, as observed in nine of the eleven cases.

\subsection{Conclusion}

Both metaheuristics prove viable for the 0/1 Knapsack Problem, each presenting distinct trade-offs. ILS consistently matched or exceeded GA solution quality while demanding a fraction of the computation time, making it the more practical choice under the current configurations. The GA's population-level diversity, while theoretically advantageous, did not translate into superior results at the chosen settings; problem-specific operators or an extended computational budget would likely alter this outcome on harder instances. The Wilcoxon test confirms that despite equal mean scores, the rank-order of outcomes differs significantly between the two algorithms, reflecting their distinct search dynamics even when they arrive at identical solution values.

% --- REFERENCES ---
\section*{References}
\addcontentsline{toc}{section}{References}

\begin{thebibliography}{9}

\bibitem{khuri1994}
Khuri, S., B\"{a}ck, T. and Heitk\"{o}tter, J., 1994, April.
The zero/one multiple knapsack problem and genetic algorithms.
In \textit{Proceedings of the 1994 ACM Symposium on Applied Computing}
(pp.\ 188--193).

\bibitem{nyathi2025}
Nyathi, T. (2025).
COS314 Artificial Intelligence Lecture Notes: Search and Optimisation.
University of Pretoria, Department of Computer Science.

\end{thebibliography}

\end{document}